\documentclass[../main.tex]{subfiles}

\graphicspath{{../image/}{../figure/}}

\begin{document}

\chapter{双曲平衡律方程组}\label{chapter:ch03-双曲平衡律方程组}

第~I 章引入的平衡律方程组所处的环境空间，在此被直观化为时空，并据此提出与定义时间方向上的双曲性这一核心概念。\cref{section:ch01-伴随平衡律}%

中讨论的通量的伴随量，在此将具体实现为熵--熵流对。

下面将给出连续介质物理中导出的双曲平衡律方程组的若干实例。

\subfile{Chapter-03/section-01.tex}

\subfile{Chapter-03/section-02.tex}

\subfile{Chapter-03/section-03.tex}

\subfile{Chapter-03/section-04.tex}

\end{document}
